%*******************************************************
% Resumen
%*******************************************************
\pdfbookmark[1]{Resumen}{Resumen}
%\addcontentsline{toc}{chapter}{\tocEntry{Resumen}}


\begingroup

\let\clearpage\relax
\let\cleardoublepage\relax
\let\cleardoublepage\relax

\chapter*{Resumen}

El análisis de datos funcionales (\acs{FDA}) es un campo de investigación estadística que proporciona un marco natural para observaciones que evolucionan sobre un dominio continuo, las cuales son cada vez más habituales en aplicaciones estadísticas. Bajo esta perspectiva, los objetos de interés son funciones aleatorias en lugar de puntos aleatorios en un espacio de dimensión finita. Debido a su creciente popularidad, se han desarrollado numerosos métodos estadísticos para datos funcionales, normalmente mediante la generalización de los enfoques existentes para los problemas multivariantes correspondientes. Sin embargo, la metodología bayesiana, aunque resulta crucial en muchos modelos de dimensión finita, sigue estando comparativamente poco desarrollada para problemas fundamentales de inferencia funcional.

En esta tesis proponemos un marco unificado de inferencia bayesiana para datos funcionales basado en espacios de Hilbert con núcleo reproductor (\acs{RKHS}), con aplicaciones a la regresión funcional y al contraste de hipótesis funcionales. Uno de los principales objetivos de esta investigación es motivar estos modelos dentro del marco funcional, a la vez que proporcionar técnicas eficientes para su aplicación en la práctica. Aunque la especificación de una distribución a priori sobre un espacio de funciones es, en general, una tarea compleja, la formulación paramétrica específica que proponemos facilita considerablemente este proceso. La distribución a posteriori resultante, que es esencial para cualquier técnica bayesiana, puede aproximarse computacionalmente mediante métodos de Montecarlo basados en cadenas de Markov de salto reversible (\acs{RJMCMC}).

Nos centramos en datos funcionales representados mediante un proceso estocástico de segundo orden \(\{X(t): t\in[0,1]\}\) con función de covarianza \(K\). Un supuesto fundamental común a todos los modelos que proponemos es que los parámetros funcionales de interés son miembros del \acs{RKHS} asociado con el núcleo \(K\). La idea central consiste en representar los parámetros funcionales como combinaciones lineales finitas de secciones del núcleo, lo que resulta en modelos de dimensión finita manejables, manteniendo al mismo tiempo una interpretación auténticamente funcional. Estos modelos pueden expresarse en términos de combinaciones lineales de las marginales del proceso \(X\), es decir, \(\sum_{j=1}^p \beta_j X(t_j)\), aprovechando la correspondencia entre el \acs{RKHS} y el espacio lineal generado por el proceso. En esta representación, el número de componentes, los coeficientes, y los instantes de tiempo asociados (a menudo denominados puntos de impacto) son cantidades desconocidas que se modelan a través de un enfoque bayesiano.

En la primera parte de la tesis desarrollamos modelos bayesianos basados en \acs{RKHS} para regresión de una variable escalar sobre una función, incluyendo tanto regresión lineal funcional como regresión logística funcional. Se asignan distribuciones a priori conjuntamente a la dimensión del modelo, los puntos de impacto y los coeficientes de regresión, y la inferencia posterior se realiza mediante \acs{RJMCMC}, permitiendo explorar modelos de distinta dimensión dentro de un único análisis bayesiano. La distribución a posteriori resultante permite considerar varias estrategias de predicción, así como un procedimiento bayesiano de selección de variables basado en la distribución marginal a posteriori de los puntos de impacto.

En la segunda parte de la tesis proponemos un contraste bayesiano de una muestra para la función de media de un proceso gaussiano, en el que se contrasta una hipótesis nula puntual frente a una alternativa funcional general mediante una aproximación \acs{RKHS}. Obtenemos una expresión cerrada de la verosimilitud mediante derivadas de Radon--Nikodym, y una distribución a priori de tipo \textit{spike-and-slab} asigna probabilidad positiva a la hipótesis nula, al tiempo que mantiene una distribución continua sobre la alternativa. La marginalización analítica de los coeficientes simplifica considerablemente el cálculo transdimensional de la distribución a posteriori, que se lleva a cabo de nuevo mediante \acs{RJMCMC}. El procedimiento resultante proporciona probabilidades a posteriori de la hipótesis nula y factores de Bayes como medidas de evidencia, así como estimaciones a posteriori y bandas creíbles para la función de media.

Además, desarrollamos una teoría asintótica para estos procedimientos bayesianos, con el fin de validar la metodología desde un punto de vista teórico. Estudiamos la consistencia de nuestros modelos cuando el tamaño muestral crece, y cuantificamos la velocidad de concentración de la distribución a posteriori. Estas garantías se ha convertido en un criterio estándar para la validación frecuentista de procedimientos bayesianos. Para la regresión funcional, establecemos la consistencia de la distribución a posteriori y las correspondientes tasas de convergencia mediante aplicaciones de resultados clásicos de Doob y Schwartz. Para el contraste funcional, se establece la consistencia tanto bajo la hipótesis nula como bajo la alternativa, con una contracción polinómica bajo la nula y exponencial bajo alternativas fijas. La teoría también se extiende al caso, relevante en la práctica, en el que la función de covarianza se estima en lugar de considerarse conocida, mostrando que estimadores de covarianza \textit{plug-in} adecuados preservan las tasas de contracción.

Nuestra contribución incluye asimismo un importante componente computacional, en forma de una implementación de muestreo mediante \acs{RJMCMC} para los modelos funcionales propuestos, que se aplica en amplios estudios numéricos. Los resultados obtenidos en estos experimentos respaldan la utilidad práctica de los métodos y muestran que nuestros procedimientos bayesianos son competitivos frente a otras alternativas frecuentistas habituales, tanto en simulaciones como en conjuntos de datos reales.

\paragraph{Palabras clave}
datos funcionales\qquad inferencia bayesiana\qquad regresión funcional\qquad contraste de hipótesis funcionales\qquad espacios de Hilbert con núcleo reproductor\qquad consistencia a posteriori\qquad Montecarlo basado en cadenas de Markov de salto reversible


\endgroup
